\scanpage{65}

Soit maintenant $A_*^{\unsure{\circ|*}}$ un \deleted{$k$-module} complexe de cochaines de $k$-modules $\frac12$ simpliciaux
\diagram{\begin{tikzcd}
0\arrow[r,dash]&A_*^0\arrow[r,dash]&A_*^1\arrow[r,dash]&A_*^2\arrow[r]&\cdots
\end{tikzcd}}
et proposons-nous de déterminer les homomorphismes de ce complexe dans
\diagram{\begin{tikzcd}
0\arrow[r,dash]&C_*^0M\arrow[r,dash]&C_*^1M\arrow[r,dash]&C_*^2M\arrow[r,dash]&\cdots,
\end{tikzcd}}
\inserted{où $M$ est un $k$-module donné}

Tout d'abord, les hom de degré $0$ (pas néc. \unsure{commutant} aux op. diff) sont les syst. d'hom
\diagram{\begin{tikzcd}
A_*^n\arrow[r,"u_n"]&C_*^nM
\end{tikzcd}}
i.e. les syst. d'hom de $k$-modules
\diagram{\begin{tikzcd}
A_n^n\arrow[r,"k_n"]&M
\end{tikzcd}}

\deleted{\unsure{On a :}}
L'hom composé
\[
d_C\circ u_n:
\]
\diagram{\begin{tikzcd}
A_*^n\arrow[r,dash]\arrow[rr,bend right=20]&C_*^nM\arrow[r]&C_*^{n+1}M
\end{tikzcd}}
correspond à l'hom composé
\diagram{\begin{tikzcd}[column sep=small]
A_{n+1}^n\arrow[rr,"k_n\partial_{n+1}"]\arrow[dr,"\partial_{n+1}"']&&M\\
&A_n^n\arrow[ur,"k_n"']&
\end{tikzcd}}

tandis que l'hom composé
\[
u_{n+1}\circ d_A:
\]
\diagram{\begin{tikzcd}
A_*^n\arrow[r]&A_*^{n+1}\arrow[r,dash]&C_*^{n+1}M
\end{tikzcd}}
correspond à l'hom composé
\diagram{\begin{tikzcd}[column sep=small]
A_{n+1}^n\arrow[rr]\arrow[dr,"d_A^n"']&&M\\
&\unsure{A_{n+1}^{n+1}}\arrow[ur,"k_{n+1}"']&
\end{tikzcd}}

\margin{64}
